\clearpage
\subsubsection{Estimate the event probability}
We will begin with a raw probability based on the YES bid/ask midpoint, recording spread, depth and timestamp. When both outcome books are usable, their midpoints can be normalised for consistency. A missing or one-sided book will be flagged, not silently replaced by an unrelated latest trade.

A simple regularised logistic model will calibrate this estimate using earlier resolved events. The wallet-informed model will then add a small number of signals: overall directional flow, the flow of historically informative wallets and market-quality controls. Coefficients will be fitted on training data, rather than chosen to produce a desired adjustment. Beta calibration is a possible robustness check if justified by the sample \cite{kull}.

\subsubsection{Select and track informative wallets}
An informative wallet is an observed account whose earlier decisions contain information beyond the contemporaneous market price. Our first score will examine whether the wallet repeatedly chose the eventual outcome at informative prices, whether prices subsequently moved in its direction, and how broad its prior history was. Returns after fees will be a supporting diagnostic; profitable liquidity provision alone is not proof of forecasting skill.

Repeated fills will be consolidated into wallet--event decisions. Balanced YES and NO holdings will be distinguished from directional exposure. Scores based on few independent events will be pulled towards the population average, a technique called shrinkage. A provisional eligibility rule is ten earlier resolved events across at least three announcement dates, with five- and twenty-event thresholds checked in sensitivity analysis. If the history cannot support this rule, the wallet signal will be unavailable rather than assigned artificial confidence.

Wallet eligibility and scores will use only records available and outcomes resolved by each forecast cut-off. New activity updates exposure; newly resolved events update later scores. Frozen snapshots prevent today's successful wallets from being retrospectively selected for earlier tests.

Contributions combine earlier skill, direction, recency and capped size so one large wallet cannot dominate. Comparing total and eligible-wallet flow tests whether selection adds information. Profiles expose these components, timestamps, concentration and disagreement; new-wallet anomalies remain separate from historical skill.

\subsubsection{Link events to assets}
The event research agent may suggest that a rate decision matters to SPY or MSFT, but numerical impacts will come from matched historical event windows and conventional market controls. The mapping must account for surprise relative to the expectation already embedded in prices. Sparse asset-specific estimates will be pooled or replaced by explicitly reviewed stress ranges.

Wallet information may improve an event probability without improving an equity-risk forecast. The event-to-asset response is a separate estimate with its own uncertainty and evidence requirements.

Each mapping will retain its event definition, horizon, sample size, estimated response range and review status. A terminal YES/NO mixture will only be used when the economic announcement falls within the risk horizon. Otherwise the MVP will offer a compatible horizon or withhold that calculation. Correlated contracts from one meeting will not be treated as independent shocks.

\clearpage
\subsubsection{Explore portfolio scenarios and tail risk}
For a selected portfolio, the engine will combine the conditional return distributions for an event occurring (YES) and not occurring (NO). Let $p$ be the selected event probability, $\mu$ a mean portfolio return and $\sigma^2$ its variance over the same horizon. The key calculations are:
\begin{align}
\mu_{\mathrm{mix}} &= p\mu_{\mathrm{YES}}+(1-p)\mu_{\mathrm{NO}}, \\
\sigma^2_{\mathrm{mix}} &= p\sigma^2_{\mathrm{YES}}+(1-p)\sigma^2_{\mathrm{NO}}
 +p(1-p)(\mu_{\mathrm{YES}}-\mu_{\mathrm{NO}})^2.
\end{align}
The last term captures uncertainty about which outcome occurs. Consequently, a higher adverse-event probability can worsen expected loss without necessarily increasing variance. For multiple holdings, we will preserve cross-asset dependence when sampling returns, then apply the user's weights to each joint scenario.

The engine will report scenario returns, horizon volatility, a chosen loss-threshold probability, 95\% VaR and 95\% expected shortfall. VaR is a threshold, not the maximum loss. Expected shortfall will be calculated as the average of the worst 5\% of the weighted loss distribution, including fractional threshold mass where needed. We will use empirical or heavy-tailed scenarios and check sensitivity to the probability and asset-response assumptions.

Paper comparisons will hold the event and model version constant while comparing the current allocation with a user-selected reduction in a supported holding and a corresponding increase in cash. Weights must sum to one and remain non-negative. Any displayed benefit must be considered alongside turnover, transaction-cost assumptions and estimation uncertainty.

\subsubsection{Volatility Forecasting \& Delta-Hedged Trading Strategy Extension}

Once the on-chain-adjusted probability $p_{\text{adj}}$ is obtained, EventLens bridges prediction-market dynamics to traditional asset volatility through a dedicated Volatility Forecasting and Delta-Hedged Trading Strategy extension. This pipeline automates contract selection, feature extraction, volatility model fitting, and delta-hedged paper trading to evaluate the forecasting performance of on-chain signals.

\paragraph{Recurring event series and contract selection}
Polymarket serves as the sole source of prediction-market data for EventLens. The platform provides continuous tracking across recurring event families, including macroeconomic releases (e.g., FOMC rate decisions, CPI releases) as well as binary Polymarket contracts on equity index closing levels (e.g., S\&P 500 / NASDAQ daily and weekly target closes) and precious metals price targets (e.g., Gold and Silver closes).

To maintain temporal consistency across continuous volatility forecasting horizons, an automated contract selection engine identifies and locks onto the closest-to-expiry market contract for any chosen event family. When event resolutions overlap or liquidity is fragmented across multiple contracts, the selection engine filters for active contracts meeting minimum volume and bid-ask spread thresholds, continuously rolling to the next active contract upon resolution.

\paragraph{Feature engineering pipeline}
The feature extraction layer constructs point-in-time predictive signals derived from the adjusted probability $p_{\text{adj}}$ alongside raw market microstructure variables. Users are given interactive control to toggle and configure exogenous inputs depending on their desired baseline.

\begin{itemize}
    \item \textbf{Probability Signal Metrics}: Includes $p_{\text{adj}}$, raw probability $p_{\text{raw}}$, probability spread $(p_{\text{adj}} - p_{\text{raw}})$, model confidence bands, and probability jump/surprise signals ($\Delta p_{\text{adj}}$ measured over historical lookback windows up to the forecast cutoff).
    \item \textbf{User-Configurable Exogenous Baselines}: Allows optional inclusion of benchmark volatility proxies (e.g., VIX, MOVE) and historical realized volatility as regime controls.
\end{itemize}

\paragraph{Realized volatility forecasting models}
The extracted signals serve as exogenous features to forecast future realized volatility over standard horizons (e.g., 1-day, 5-day). The core target over $H$ trading days is defined as the sum of squared daily log returns:
\begin{equation}
RV_{t,H}=\sum_{h=1}^{H} r_{t+h}^{2}.
\end{equation}

EventLens implements three core volatility model structures:
\begin{itemize}
    \item \textbf{HAR-RV-X (Heterogeneous Autoregressive with Exogenous Signals)}: Extends the classical HAR-RV framework (which models realized volatility via daily, weekly, and monthly components) by adding linear exogenous terms for $p_{\text{adj}}$ and probability jump signals. Assumes a linear, multi-frequency structure for volatility persistence.
    \item \textbf{GARCH-X}: Incorporates exogenous prediction-market jump and log-odds signals directly into the conditional variance equation of a standard GARCH model. Assumes stationarity and clustering in conditional volatility returns.
    \item \textbf{Machine Learning Baselines (XGBoost / Random Forest)}: Non-linear tree ensemble models fitted on historical volatility features, probability spreads, and regime signals. Evaluated using standard robust loss functions such as QLIKE (Quasi-Likelihood) and Mean Squared Error (MSE) to prevent extreme volatility outliers from distorting model fitting.
\end{itemize}

\paragraph{Delta-hedged backtesting framework}
To demonstrate the practical utility of the forecasted realized volatility, EventLens runs an automated delta-hedged backtest. The backtest engine compares the model's forecasted realized volatility against current option-implied volatility (IV) from option market quotes to identify implied-versus-forecasted variance gaps. 

When a variance gap is detected, the strategy simulates a delta-neutral position (e.g., ATM options hedged continuously against the underlying asset) to harvest the variance risk premium. Costs, bid/ask spreads, hedge frequency, and execution assumptions are included. This simplified backtest framework provides the user with an intuitive evaluation of forecast performance, displaying backtested Sharpe ratios, maximum drawdowns, and simulated trade logs directly on the research dashboard.